\section{De ROS a ROS~2}

\subsection{Origen de ROS}

ROS (Robot Operating System) nació en 2007 en el laboratorio de inteligencia artificial de Stanford bajo el proyecto STAIR (STanford Artificial Intelligence Robot). Su objetivo original era resolver un problema concreto: los equipos de robótica repetían los mismos componentes de software en cada proyecto nuevo porque no había una plataforma común sobre la que construir~\cite{quigley2009ros}. Willow Garage, una empresa de robótica fundada en 2006 en Menlo Park, adoptó el proyecto en 2008 y lo convirtió en el estándar de facto de la investigación robótica.

La propuesta de ROS era elegante: en lugar de definir una arquitectura monolítica, definió un conjunto de convenciones. Los procesos (nodos) se comunicarían mediante mensajes tipados sobre canales nombrados (tópicos), sin saber nada el uno del otro. El código podía reutilizarse entre robots completamente distintos porque la interfaz era siempre la misma: publicar y suscribirse a tópicos estándar. Un nodo de teleoperación escrito para un robot podía controllar cualquier otro robot que publicara en \texttt{/cmd\_vel} con el tipo \texttt{geometry\_msgs/Twist}.

El resultado fue una comunidad de contribución masiva. Para 2013 existían miles de paquetes ROS~1 que cubrían desde percepción visual hasta planificación de trayectorias. El robot PR2 de Willow Garage fue el catalizador: muchos laboratorios lo adoptaron y construyeron sobre el mismo stack de software, creando un ecosistema compartido sin precedente en robótica~\cite{quigley2009ros}.

\subsection{Limitaciones de ROS~1}

El éxito de ROS~1 expuso sus limitaciones de diseño. ROS~1 fue construido con supuestos que funcionaban perfectamente en laboratorio pero que fallaban en producción~\cite{maruyama2016exploring}:

\textbf{Punto único de fallo: el rosmaster.} Toda comunicación en ROS~1 depende de un proceso central llamado \texttt{rosmaster}. Cuando un nodo quiere publicar en un tópico, registra su intención con el master. Cuando otro nodo quiere suscribirse, consulta al master por la dirección del publicador. El master es el directorio central del sistema. Si ese proceso falla, se desconecta, o la red entre él y algún nodo se interrumpe, el sistema entero queda inutilizable. Esto hace imposible construir sistemas verdaderamente distribuidos en múltiples máquinas sin un punto central de fallo.

\textbf{Sin garantías de comunicación.} El protocolo de transporte de ROS~1 (TCPROS y UDPROS) no ofrece ninguna garantía de calidad de servicio. Los mensajes se entregan si el suscriptor está activo; si no, se pierden. No hay mecanismo de historial, persistencia, ni reintentos. Para un stream de cámara esto es aceptable, pero para comandos críticos es un riesgo real.

\textbf{Python~2 y C++03.} ROS~1 fue construido sobre Python~2, cuyo fin de vida oficial fue en enero de 2020. Mantener código ROS~1 después de esa fecha implica usar un intérprete sin soporte de seguridad. Las APIs de C++ también usaban estándares antiguos, limitando el uso de características modernas del lenguaje.

\textbf{Sin soporte de tiempo real.} El diseño de ROS~1 no contempló nunca requisitos de tiempo real. Los tiempos de entrega de mensajes no están garantizados, lo que impide su uso en lazos de control con restricciones temporales estrictas.

\textbf{Sin seguridad.} Toda la comunicación en ROS~1 es sin cifrado y sin autenticación. Cualquier nodo en la red puede suscribirse a cualquier tópico, inyectar mensajes, o incluso llamar a servicios. Esto es inaceptable en sistemas con acceso a internet o en entornos industriales.

\subsection{El rediseño: ROS~2}

ROS~2 no es una versión más de ROS~1. Es un rediseño completo que mantiene los conceptos de alto nivel (nodos, tópicos, servicios, parámetros) pero reemplaza toda la infraestructura subyacente~\cite{macenski2022robot}. El diseño fue guiado por cinco objetivos explícitos:

\textbf{Eliminación del master.} ROS~2 delega el descubrimiento de nodos al middleware DDS (Data Distribution Service), un estándar de la industria que implementa descubrimiento distribuido sin servidor central. Los nodos se anuncian entre sí directamente en la red usando el protocolo RTPS (Real-Time Publish Subscribe). La desaparición de un nodo no afecta a ningún otro.

\textbf{Calidad de servicio configurable.} Heredando las políticas QoS de DDS, cada publicador y suscriptor puede configurarse independientemente con distintos niveles de confiabilidad, historial y durabilidad. Esto permite que el mismo sistema use comunicación confiable para comandos de seguridad y comunicación de mejor esfuerzo para streams de alta frecuencia.

\textbf{Soporte de tiempo real.} La arquitectura de ROS~2 fue diseñada explícitamente para ser compatible con ejecutores de tiempo real. La capa de abstracción del middleware (RMW) puede ser reemplazada por implementaciones deterministas. La configuración completa requiere un kernel Linux con parches PREEMPT\_RT y una implementación DDS que ofrezca garantías temporales.

\textbf{Seguridad nativa.} SROS~2 (Secure ROS~2) es la capa de seguridad integrada en ROS~2. Provee cifrado de comunicaciones, autenticación de nodos mediante certificados, y control de acceso por tópico. Se activa sin cambiar el código de aplicación.

\textbf{Python~3 y C++ moderno.} ROS~2 fue construido desde cero sobre Python~3 y C++14/17. Las APIs de \texttt{rclpy} y \texttt{rclcpp} usan características modernas: type hints, lambdas, smart pointers, y estructuras de datos estándar.

La arquitectura interna también cambió radicalmente. En ROS~1, la biblioteca de cliente era una sola capa monolítica. En ROS~2, la pila se divide en tres niveles: \texttt{rcl} (ROS Client Library), una biblioteca en C puro que implementa la lógica de ROS independiente del lenguaje; \texttt{rmw} (ROS Middleware), la capa de abstracción que oculta la implementación DDS concreta; y \texttt{rclpy}/\texttt{rclcpp}, los bindings de alto nivel para Python y C++ respectivamente. Esta separación permite cambiar la implementación DDS sin recompilar el código de la aplicación.

\subsection{Coexistencia: el puente ros1\_bridge}

La transición de ROS~1 a ROS~2 no fue inmediata. Para facilitar la migración gradual, existe el paquete \texttt{ros1\_bridge}~\cite{ros1bridge}, que actúa como traductor bidireccional entre los dos sistemas. El puente corre como un nodo simultáneamente registrado en ambos entornos: suscribe tópicos de ROS~1 y los republica en ROS~2, y viceversa.

Su uso requiere tener ambas distribuciones instaladas simultáneamente y sourcear ambos entornos en la misma terminal. Es una solución de migración, no una arquitectura de producción: la traducción introduce latencia adicional y no todos los tipos de mensajes tienen mapeo automático. Los proyectos nuevos no deben diseñarse con ros1\_bridge como componente permanente.
